\newpage
\begin{adjustwidth}{-5pt}{-2pt}
\begin{singlespace}
\tcbset{colback=white!10!white}
\begin{tcolorbox}[title=\textbf{\section{Model Card - \seamless}},
    breakable, sharp corners, boxrule=0.5pt] 
\begin{mcsection}{Model Details \footnote{For this card, we use the template from \cite{10.1145/3287560.3287596}. }} 
\item Person or organization developing model: \textit{Developed by FAIR, Meta}
\item Model date: \textit{November 30, 2023}
\item Model version: \seamless
\item Model type: \seamlessstreaming translation model with two PRETSSEL acoustic models (6 and 36 languages) and two HiFi-GAN mel-vocoder (16kHz and 24kHz) \
\begin{itemize}
    \item  \it License: custom research license
\item Where to send questions or comments about the model: \\
\url{https://github.com/facebookresearch/seamless_communication/issues}
\end{itemize} 
\end{mcsection}


\begin{mcsection}{Intended Use}
    \item Primary intended uses: \textit{\seamless model is an expressive multilingual streaming translation model. It allows for:
    \begin{itemize}
        \item Simultaneous translation from 100 source languages in speech into a selection of 6 or 36 target languages in speech.
        \item Preservation of sentence-level prosody and vocal style.
    \end{itemize}
Information on how to use the model can be found in seamless\_communication repository}
\item Primary intended users: \textit{Primary users are researchers and machine translation (speech and text) research community. }
\item Out-of-scope use cases: \textit{\seamless model is a research model and is not released for production deployment. 
} 
\end{mcsection}

\begin{mcsection}{Metrics}
    \item Quality: \textit{ASR-BLEU}, \textit{vocal style similarity}, \comparator

    \item Latency: \textit{Ending offset.}
\end{mcsection}


\begin{mcsection}{Evaluation Data}
    \item Datasets: \textit{\fleurs}
\item Motivation: \textit{We used \fleurs as it provides
an n-way parallel speech and text dataset in 102 languages, on which we can evaluate \seamless models on multiple tasks.}

\end{mcsection}

\begin{mcsection}{Training Data}
\item We used parallel multilingual speech from a variety of sources to train models.
\end{mcsection}



\begin{mcsection}{Ethical Considerations}
 \item \textit{In this work, we took a comprehensive approach to prioritize human users and minimize risks that could be transferred to them. While we have documented various evaluation and responsible AI techniques deployed in our work, here are some additional points to highlight. For one, many languages chosen for this study are low-resource languages.
 While quality translation could improve world readiness and information access for many in these communities, such access could also make groups with lower levels of digital literacy more vulnerable to misinformation or online scams. The latter scenarios could arise if bad actors misappropriate our work for nefarious activities, which we conceive as an example of unintended use.
 Finally, although we did our best to optimize for translation quality, toxic, biased, or false outputs produced by the model could remain. These could have an adverse impact on those who rely on these translations to make important decisions (particularly when related to health and safety).} 
\end{mcsection}


\begin{mcsection}{Caveats and Recommendations}
  \item Limitations: \textit{Researchers should consider implementing additional integrity mitigations for ``added toxicity'' when using the model in a research application.}
\end{mcsection}


\end{tcolorbox}
\end{singlespace}
\end{adjustwidth}
